\par\noindent\textbf{(1)} 在体积 \(V\) 内，两偏振模式每单位角频率状态密度 \(g(\omega)/V=\omega^2/(\pi^2c^3)\)。每模平均能量 \(\hbar\omega/(e^{\beta\hbar\omega}-1)\)，故 \(u_\omega=\hbar\omega^3/[\pi^2c^3(e^{\hbar\omega/(k_BT)}-1)]\)，零点能不贡献热辐射通量。
\par\noindent\textbf{(2)} 令 \(x=\hbar\omega/(k_BT)\)，使用 \(\int_0^\infty x^3/(e^x-1)dx=\pi^4/15\)，得 \(u=\pi^2(k_BT)^4/(15\hbar^3c^3)\)。等方辐射穿过平面半空间的能流 \(j^*=cu/4=\sigma T^4\)，\(\sigma=\pi^2k_B^4/(60\hbar^3c^2)\)。
\par\noindent\textbf{(3)} 按波长频谱 \(u_\lambda=(8\pi hc/\lambda^5)/(e^{hc/(\lambda k_BT)}-1)\)。取 \(x=hc/(\lambda k_BT)\)，驻点满足 \(5(1-e^{-x})=x\)，非零正根 \(x\approx4.9651\)。所以 \(\lambda_{\max}T=hc/(4.9651k_B)\)。
\par\noindent\textbf{(4)} 谱密度变换必须乘Jacobian：\(u_\lambda=u_\nu|d\nu/d\lambda|\propto u_\nu c/\lambda^2\)。因此频率谱极大值满足 \(3(1-e^{-x})=x\)，与波长谱的 \(5(1-e^{-x})=x\) 不同。不能仅令 \(\lambda=c/\nu_{\max}\) 求波长峰。
